% ПЗ-3 НИТУ МИСИС — статья по утверждённому классу преподавателя.
% Сборка: tectonic article.tex; проект использует оригинальный misisarticle.cls.
% Экспериментальные методы, не подтверждённые в НИР, описаны как планируемые.
% Файл misisarticle.cls не заменять: сохраняется оригинал в форке.
\documentclass{misisarticle}
\usepackage{amsthm,mathtools,tabularx,tikz}
\usetikzlibrary{arrows.meta,positioning,shapes.geometric,fit,calc}
\newcommand{\eng}[1]{\foreignlanguage{english}{#1}}
\newcommand{\code}[1]{\texttt{#1}}
\newcommand{\Ind}{\mathbf{1}}
\DeclareMathOperator*{\argmin}{arg\,min}
\DeclareMathOperator*{\argmax}{arg\,max}
\newtheorem{proposition}{Утверждение}
\title{Исследование распределённого инференса Transformer-моделей с послойным шардированием на базе Ray}
\misisauthor{Соколов Филипп Олегович}{БПМ-23-ПО-3, кафедра инженерной кибернетики}{НИТУ МИСИС, Москва}{}
\grnti{50.39.19}
\asjc{1712}
\oecd{1.2}
\abstract{Рассматривается задача оценки и сокращения коммуникационных накладных расходов при распределённом авторегрессионном инференсе Transformer-моделей в условиях ограниченных вычислительных ресурсов. Исходной основой является реализованный CPU-прототип с послойным размещением компонентов GPT-2-совместимой модели на трёх логических узлах Ray, развёрнутых в контейнерах одного физического компьютера. Формализованы четыре варианта обработки выходных логитов: передача полного тензора, передача последней позиции, вычисление логитов только последней позиции и локальный выбор токена завершающим исполнителем. Выведена аналитическая оценка полезного объёма данных и предложена методика раздельной оценки вычислительных и коммуникационных эффектов. В качестве следующего исследовательского этапа сформулирована задача выбора границ послойного разбиения с учётом стоимости вычислений, обмена и ограничений памяти; предложено решение методом динамического программирования для аддитивной модели задержки. Спроектирован авторский набор из 144 нагрузочных сценариев, предусмотрены контроль корректности генерации и сравнение с монолитным исполнением. В статье представлены постановка и методы; модификации, оптимизационный алгоритм и расширенные измерения описаны как планируемые, без утверждений о достигнутом ускорении.}
\keywords{распределённый инференс; послойное шардирование; Ray Actors; авторегрессионная генерация; передача промежуточных тензоров; коммуникационные накладные расходы}
\begin{document}
\maketitle
\section{Введение}

Инференс больших языковых моделей становится одной из важных вычислительных задач: размер весов, длина контекста и число одновременных запросов предъявляют различные требования к памяти, вычислительным ресурсам и обмену данными. Если модель нельзя эффективно выполнить на одном устройстве, применяют модельный параллелизм, распределяя разные части её вычислительного графа между исполнителями. При послойном разбиении последовательные группы блоков Transformer закрепляются за разными узлами, а промежуточные представления передаются между ними. Такой подход необходимо отличать от параллельной обработки запросов несколькими полными копиями модели \cite{petals2023,deepspeed2022}.

Распределение вычислений само по себе не гарантирует ускорения. Передача активаций, сериализация тензоров, удалённые вызовы, хранение весов и недостаточная загрузка этапов могут перевесить потенциальную выгоду. Исследования последних лет показывают существенную роль коммуникационных механизмов и баланса этапов даже в высокопроизводительных системах на нескольких GPU \cite{jekel2026,hwang2026}. В малых CPU-конфигурациях влияние накладных расходов может быть особенно заметно. При этом мощный GPU-кластер не является необходимым условием для исследования \emph{корректности алгоритма, стоимости информационного обмена и предсказательной модели размещения слоёв}.

В качестве отправной точки используется ранее созданный автором учебно-экспериментальный прототип распределённого инференса \code{sshleifer/tiny-gpt2}, реализованный с помощью PyTorch, Hugging Face Transformers и Ray Actors. В нём модель разделена на три функциональные части и последовательно исполняется в трёх контейнерах Docker Compose на \emph{одной физической машине}. Прототип подтверждён запуском Ray-кластера, проверкой размещения исполнителей, генерацией текста и автоматическими тестами. Однако сама tiny-модель имеет лишь два Transformer-блока и служит проверке архитектуры, а не оценке промышленного многоузлового инференса.

Наиболее конкретное обнаруженное ограничение исходной реализации связано с возвратом полных выходных логитов. При жадном декодировании для продолжения последовательности требуется выбрать один следующий токен по логитам \emph{последней позиции}. Возврат логитов для всех позиций создаёт избыточный объём данных. Но важно установить, какая доля возможной экономии связана с сокращением \emph{передачи}, а какая~--- с уменьшением \emph{вычислений} при формировании логитов. Поэтому простое сравнение «исходная реализация~--- перенос \code{argmax}» недостаточно для причинно интерпретируемого эксперимента.

Другой перспективный аспект~--- выбор границ разделения модели. Для реальных многослойных моделей число возможных разбиений быстро растёт. При неоднородной скорости узлов, разных объёмах доступной памяти и стоимости передачи данных равномерное деление по числу слоёв не обязательно минимизирует задержку. В то же время для одного последовательного запроса нельзя без дополнительных допущений подменять минимизацию \emph{суммарной задержки} минимизацией времени самого медленного этапа; последняя задача непосредственно связана с производительностью перекрывающегося многозапросного конвейера \cite{hwang2026}.

\textbf{Цель исследования}~--- разработать воспроизводимую методику анализа коммуникационных затрат послойного инференса и сформулировать ресурсно-ограниченный способ выбора размещения компонентов модели. В рамках текущего этапа решаются задачи: (1) формализация объёмов обмена и режимов обработки логитов; (2) постановка гипотез о влиянии изменений на задержку; (3) формулировка задачи разбиения с ограничениями памяти; (4) проектирование алгоритма и контрольных проверок; (5) описание авторских нагрузочных данных и реалистичного эксперимента на доступном оборудовании.

\textbf{Границы выполненного и планируемого.} Исходный вариант с полным возвратом логитов уже реализован в НИР. Сокращённые режимы, алгоритм выбора границ разбиения и серия сравнительных экспериментов в настоящей статье предлагаются к реализации. Результаты ускорения, оптимальность на реальном оборудовании и показатели многомашинного кластера не объявляются полученными.

\section{Анализ существующих исследований}

\subsection{Распределение модели и ограничения памяти}

Система \emph{Petals} \cite{petals2023} позволяет исполнять последовательные блоки больших языковых моделей на узлах нескольких участников. Она демонстрирует практически значимый сценарий распределения слоёв по сети, но включает задачи динамической маршрутизации, управления участниками и отказоустойчивости. В предлагаемой работе изучается более контролируемый стенд: исполнители фиксированы, а основной предмет анализа~--- стоимость обмена между стадиями и воспроизводимость измерений.

\emph{DeepSpeed Inference} \cite{deepspeed2022} объединяет методы параллелизма и размещения вычислений для крупных Transformer-моделей в многоускорительной среде. Исследование показывает, что оптимальная организация инференса зависит от ограничивающего ресурса и выбранной нагрузки. Непосредственное сравнение скорости DeepSpeed с локальным CPU-прототипом не имеет смысла из-за разницы масштабов, однако принципы моделирования ограничений являются релевантными.

\emph{FlexGen} \cite{flexgen2023} ориентирован на генерацию при ограниченной GPU-памяти. В нём совместно используются GPU, RAM и накопитель, а выбор стратегии размещения данных формулируется как задача оптимизации. В отличие от FlexGen, здесь предполагается оптимизировать границы \emph{последовательных групп Transformer-слоёв} и анализировать стоимость межпроцессного обмена в Ray. Идея учёта ресурсоёмкости отдельных компонентов при этом концептуально сходна.

\subsection{Обслуживание запросов, коммуникации и балансировка}

В системе \emph{vLLM} \cite{vllm2023} механизм PagedAttention применяется для управления KV-кэшем при множестве запросов. Система \emph{DistServe} \cite{distserve2024} исследует раздельное размещение фаз обработки входного контекста и декодирования. Обе работы показывают значимость формы нагрузки и состояния вычислительного процесса; текущая версия нашего прототипа не содержит KV-кэш и многозапросную конвейеризацию.

Jekel и Bhatele \cite{jekel2026} исследуют коммуникационные ограничения многоузлового GPU-инференса, уделяя особое внимание коллективным операциям при тензорном параллелизме. Их задача отличается от последовательного обмена активациями между Ray Actors, поэтому численные результаты нельзя переносить на локальную CPU-систему. Тем не менее работа обосновывает необходимость раздельного анализа вычислительных и коммуникационных компонентов задержки.

Hwang и Ahn \cite{hwang2026} исследуют балансировку стадий конвейерного обслуживания LLM и показывают, что дисбаланс и простои способны ограничивать пропускную способность. Это подтверждает значимость оптимизации размещения слоёв, но подчёркивает необходимость различать два режима: \emph{последовательную обработку одного запроса} и \emph{перекрывающийся конвейер нескольких запросов}.

Фреймворк \emph{Ray} \cite{ray2018} обеспечивает удалённые объекты и планирование по логическим ресурсам. Выбор Ray является инженерным решением для прототипа, а не предметом заявляемой научной новизны. Логические ресурсы Ray не обеспечивают жёсткой привязки к отдельным физическим ядрам \cite{rayresources}.

\begin{table*}[t]
\centering\small
\caption{Сопоставление работ и исследовательской постановки}
\label{tab:related}
\begin{tabularx}{\textwidth}{@{}p{2.55cm}X X@{}}
\toprule
\textbf{Подход} & \textbf{Основной предмет} & \textbf{Связь с данной работой}\\
\midrule
Petals \cite{petals2023} & Размещение блоков LLM по сети & Послойный обмен; здесь фиксированные исполнители\\
DeepSpeed \cite{deepspeed2022} & Инференс крупных моделей на GPU & Учёт вычислительных и ресурсных ограничений\\
FlexGen \cite{flexgen2023} & Размещение данных при дефиците памяти & Оптимизационная постановка при ограниченных ресурсах\\
vLLM, DistServe \cite{vllm2023,distserve2024} & KV-кэш и организация обслуживания & Ориентиры для будущего расширения; нет в исходной системе\\
Jekel--Bhatele \cite{jekel2026} & Коммуникации многоузлового GPU-инференса & Обоснование декомпозиции накладных расходов\\
Hwang--Ahn \cite{hwang2026} & Балансировка конвейерных стадий & Различие между задержкой и конвейерной пропускной способностью\\
Настоящая работа & Сериализация выходных данных и выбор разбиения & Воспроизводимая CPU-постановка с проверяемыми ограничениями\\
\bottomrule
\end{tabularx}
\end{table*}

Обзор выявляет конкретную нишу исследования: на простом управляемом стенде последовательно выделить эффект объёма возвращаемых логитов, стоимость их вычисления и влияние выбора границ разбиения на прогнозируемую задержку. Такой подход не заменяет промышленного сервинга, но позволяет получить воспроизводимые результаты о свойствах исследуемой архитектуры.

\section{Постановка задачи}

\subsection{Модель вычислительного графа и режимы выполнения}

Пусть задана авторегрессионная decoder-only Transformer-модель с $N$ блоками, скрытой размерностью $H$, словарём из $V$ токенов и параметрами $\theta$. Для исходной последовательности $x_{1:L_0}$ на шаге $t$ вычисляется распределение следующего токена, а при жадном декодировании выбирается
\begin{equation}
 x_{L_t+1}=\argmax_{v\in\{1,\ldots,V\}} z_{L_t,v},\qquad L_t=L_0+t-1,
 \label{eq:greedy}
\end{equation}
где $z_{L_t,v}$~--- логит токена $v$ в последней позиции. Для сопоставимости режимов фиксируются веса, токенизатор, тип данных, входы и детерминированное правило декодирования.

Пусть модель разделена на $K$ \emph{упорядоченных} исполнителей. Границы разделения $0=b_0\le b_1\le\ldots\le b_K=N$ задают интервалы $[b_{i-1},b_i)$ Transformer-блоков. На первом исполнителе дополнительно размещены эмбеддинги, на последнем~--- конечная нормализация и выходная голова. При $N<K$ допускаются пустые диапазоны блоков, если за соответствующим исполнителем остаются необходимые вычисления. Для содержательного сравнения распределения многослойной модели в дальнейшем предполагается $N\ge K$ и непустые диапазоны.

Исходная реализация на Ray относится к \emph{последовательному исполнению одного запроса}. Она не реализует перекрытие вычислений разных микропакетов, поэтому в рассматриваемом локальном варианте первичной метрикой является \emph{полная задержка} одного запроса, а не загрузка самого медленного этапа.

\subsection{Контрольные и модифицированные режимы}

Вводятся пять режимов (табл.~\ref{tab:variants}). Такое разделение позволяет не смешивать уменьшение объёма выходных данных с уменьшением объёма вычислений в выходной голове.

\begin{table*}[t]
\centering\small
\caption{Планируемая матрица сравнительных реализаций}
\label{tab:variants}
\begin{tabularx}{\textwidth}{@{}p{1.15cm}X p{4.0cm}@{}}
\toprule
\textbf{Код} & \textbf{Что выполняется} & \textbf{Статус}\\
\midrule
B0 & Монолитная модель; локальный выбор токена & Контрольный монолитный путь в исходном проекте\\
D0 & Распределённая модель, возврат $[B,L,V]$ логитов & \textbf{Реализовано} в исходной НИР\\
D1 & Полные логиты вычисляются, но возвращается $[B,V]$ & Планируется\\
D2 & Только логиты последней позиции вычисляются и возвращаются & Планируется\\
D3 & Логиты последней позиции и \code{argmax} на последнем исполнителе, возврат ID & Планируется\\
\bottomrule
\end{tabularx}
\end{table*}

\textbf{Гипотеза $H_1$:} последовательное устранение избыточных выходных данных в D1--D3 уменьшает полезный объём обмена относительно D0. \textbf{Гипотеза $H_2$:} сокращение вычислений и передачи может уменьшать общую задержку, но знак и величина эффекта зависят от среды, длины контекста и накладных расходов Ray. \textbf{Гипотеза $H_3$:} для неоднородных исполнителей и неравномерной стоимости слоёв модель стоимости позволяет найти допустимое разбиение с меньшей \emph{прогнозной} задержкой, чем фиксированное равномерное разбиение. $H_3$ не утверждает автоматического ускорения на единственном физическом CPU.

\subsection{Критерии корректности и ограничения}

Для каждого входа необходимо сверять последовательности \emph{идентификаторов} генерируемых токенов; различия не удаляются из результатов, а протоколируются. При проверке внутренних тензоров допускается числовой допуск, обоснованный реализацией операций и точностью вычислений. Случаи изменения первого токена, ошибки маршрутизации, переполнения памяти или рассинхронизации времени должны сохраняться как ошибки эксперимента.

В качестве показателей определяются: (i) $t_{\mathrm{gen}}$, полная задержка; (ii) $\mathrm{TPS}=T/t_{\mathrm{gen}}$, скорость выходных токенов для фиксированного одиночного запроса; (iii) $V_{\mathrm{payload}}$, аналитический полезный объём; (iv) длительность вычисления выходной головы и удалённых вызовов; (v) при наличии измерителей~--- реальный транспортный трафик и пиковая резидентная память процессов. TPS одиночного запроса не следует отождествлять с пропускной способностью многозапросного сервера.

Основная экспериментальная площадка~--- один компьютер с несколькими логическими узлами Ray. Полученные на нём результаты не могут сами по себе подтверждать масштабирование на физически независимых серверах или скоростях Ethernet.

\section{Методы и организация исследования}

\subsection{Подтверждённая архитектура исходного прототипа}

Предыдущая НИР реализована с помощью PyTorch, Hugging Face Transformers, Ray Actors, FastAPI и Docker Compose. В качестве небольшой контрольной модели используется \code{sshleifer/tiny-gpt2}: $N=2$ Transformer-блока, $H=2$, $V=50\,257$. В рамках одного физического компьютера запускаются контейнеры \code{ray-head}, \code{ray-worker-1}, \code{ray-worker-2}, зарегистрированные как три логических Ray-узла.

Исполнитель 0 содержит эмбеддинги и блок $[0,1)$, исполнитель 1~--- блок $[1,2)$, исполнитель 2~--- заключительную нормализацию и \code{lm\_head} при пустом диапазоне $[2,2)$. На каждом шаге генерации через удалённые вызовы последовательно проходят промежуточные тензоры. В исходном варианте отсутствует KV-кэш: текущая последовательность обрабатывается целиком на каждом шаге. Это упрощает контроль корректности, хотя и увеличивает вычислительные затраты при росте длины контекста.

\begin{figure*}[t]
\centering
\begin{tikzpicture}[every node/.style={font=\footnotesize},node distance=5mm,
box/.style={draw,rounded corners,align=center,minimum height=9mm,text width=5.0cm,fill=gray!6},
arr/.style={-{Stealth[length=2mm]},line width=0.65pt}]
\node[box] (a) {Ray head / shard 0\\эмбеддинги + Transformer-блок 0};
\node[box,below=of a] (b) {Ray worker 1 / shard 1\\Transformer-блок 1};
\node[box,below=of b] (c) {Ray worker 2 / shard 2\\финальная нормализация + lm\_head};
\node[draw,rounded corners,align=center,text width=3.0cm,right=18mm of b,minimum height=16mm] (d) {Управляющий процесс\\выбор / добавление токена};
\draw[arr] (a)--node[right]{hidden states}(b);
\draw[arr] (b)--node[right]{hidden states}(c);
\draw[arr] (c.east)-|node[pos=0.24,right,font=\scriptsize]{logits / token\_id}(d.south);
\draw[arr,dashed] (d.north)|- (a.east);
\node[font=\scriptsize,align=center,below=6mm of c] {Три контейнера Ray на \textbf{одном} физическом компьютере};
\end{tikzpicture}
\caption{Исполнение шагов генерации в исходной архитектуре; режимы D1--D3 изменяют только обработку выходных логитов.}
\label{fig:pipeline}
\end{figure*}

По материалам выполненной НИР на таком стенде подтверждены три активных Ray-узла и успешная генерация 20 токенов. При этом 24 CPU, которые отображались как суммарные ресурсы Ray, являются \emph{логическими объявлениями ресурсов}, а не 24 независимыми физическими ядрами \cite{rayresources}. При переносе на MacBook Pro с Apple M1 Pro и 16~ГБ памяти параметры Docker и фактическая загрузка должны быть проверены повторно: прежние значения скорости нельзя автоматически распространять на новый компьютер.

\subsection{Модель коммуникационных расходов}

Пусть $B$~--- размер пакета, $d_z$~--- число байт на один логит, $T$~--- число генерируемых токенов, $L_t=L_0+t-1$~--- длина контекста на шаге $t$. Для режима D0 сырой объём \emph{полезных выходных данных} от завершающего исполнителя составляет
\begin{equation}
 V_{\mathrm{D0}} = B V d_z\sum_{t=1}^{T}L_t
 = BVd_z\left(TL_0+\frac{T(T-1)}{2}\right).
 \label{eq:full}
\end{equation}
Для D1 и D2 выходной объём одинаков, хотя стоимость вычислений различается:
\begin{equation}
 V_{\mathrm{D1}}=V_{\mathrm{D2}}=B V d_z T.
 \label{eq:last}
\end{equation}
Для D3 выход состоит из идентификаторов токенов длиной $d_{\mathrm{id}}$ байт:
\begin{equation}
 V_{\mathrm{D3}}=B d_{\mathrm{id}} T.
 \label{eq:id}
\end{equation}

Для выполненного ранее контрольного сценария $B=1$, $L_0=5$, $T=20$, $V=50\,257$ и $d_z=4$ получаем $\sum L_t=290$. Тогда
\begin{equation}
 V_{\mathrm{D0}}=58\,298\,120\ \text{байт}\approx55{,}60\ \text{МиБ},
 \label{eq:example}
\end{equation}
а для D1--D2~--- $4\,020\,560$ байт (примерно $3{,}83$~МиБ). При передаче идентификаторов \code{int64} в D3~--- $160$ байт. Все три величины относятся \emph{только к выходному направлению} и являются расчётными, а не замерами транспортного трафика.

Для промежуточных hidden states на одной границе при последовательной передаче полного контекста величина составляет $Bd_hH\sum_t L_t$, где $d_h$~--- размер элемента активации. В исходной НИР суммарный объём двух границ оценён в $4\,640$ байт, а отдельно передаваемых масок~--- ещё в $4\,640$ байт. Поэтому минимальный сырой объём всех рассматриваемых тензоров в D0 равен $58\,307\,400$ байт (около $55{,}61$ МиБ) без заголовков Ray, сериализации и дополнительных сообщений.

Модель задержки должна учитывать как минимум
\begin{equation}
 t_{\mathrm{gen}}\approx t_{\mathrm{compute}}+t_{\mathrm{rpc}}
 +t_{\mathrm{serialize}}+t_{\mathrm{transport}}+t_{\mathrm{schedule}}.
 \label{eq:latency}
\end{equation}
Эта декомпозиция является \emph{моделью для измерений}, а не установленным экспериментально равенством: компоненты могут перекрываться. Уменьшение $V_{\mathrm{payload}}$ не гарантирует пропорционального уменьшения $t_{\mathrm{gen}}$.

\subsection{Разделение вычислительного и коммуникационного эффекта}

Предлагается последовательная абляция D0--D3. При переходе D0$\to$D1 сокращается объём возвращаемого сообщения без обязательного сокращения вычисления выходной головы. При D1$\to$D2 логиты вычисляются только для последней позиции: для причинной GPT-2-модели это допускается послойной нормализацией и независимым применением выходной линейной головы к каждой позиции. При D2$\to$D3 выбор \code{argmax} переносится внутрь последнего Ray Actor, и вместо вектора возвращается один ID.

Каждый переход проверяется на идентичность последовательностей токенов эталонному B0 при зафиксированных весах и численной конфигурации. Вклад отдельных изменений оценивается попарными сравнениями D0--D1, D1--D2, D2--D3. Такое разложение позволяет не приписывать экономии передачи данных ускорение, фактически достигнутое сокращением вычислений.

\subsection{Математическая задача выбора послойного разбиения}

Для дальнейшей работы с многослойной GPT-2-совместимой моделью вводится отдельная оптимизационная постановка. Пусть $N\ge K$, блоки сохраняют исходный порядок, а $\omega$ обозначает один нагрузочный сценарий (длину входа, генерации и размер пакета). Для исполнителя $i$ определяются оценённое время вычислений $C_i(p,q;\omega)$ для блоков $[p,q)$, потребность в памяти $M_i(p,q;\omega)$ и доступный предел $M_i^{\max}$. Эмбеддинги относятся к первому, а выходная голова~--- к последнему исполнителю и учитываются в их функциях стоимости.

Для последовательного исполнения вводится аддитивная \emph{прогнозная} стоимость
\begin{equation}
 \widehat T(\mathbf b;\omega)
 =\sum_{i=1}^{K} C_i(b_{i-1},b_i;\omega)
 +\sum_{i=1}^{K-1} Q_i(b_i;\omega),
 \label{eq:cost}
\end{equation}
где $Q_i$~--- стоимость передачи активаций на границе между исполнителями, в том числе задержка вызова, сериализация и оценка передачи данных. Для \emph{независимых} физических каналов возможна приближённая форма $Q_i\approx\alpha_i+S_i/\beta_i$, где $S_i$~--- число байт, $\beta_i$~--- эффективная скорость линии; для локальных контейнеров параметры следует измерять как межпроцессные затраты, а не выдавать за Ethernet-характеристики.

На конечном наборе нагрузок $\Omega$ с весами $w_\omega\ge0$ и $\sum_\omega w_\omega=1$ оптимизационная задача имеет вид
\begin{align}
 \min_{0=b_0<b_1<\cdots<b_K=N}\;&\sum_{\omega\in\Omega}w_\omega\widehat T(\mathbf b;\omega),\label{eq:opt}\\[-2mm]
 \text{при условиях}\quad &M_i(b_{i-1},b_i;\omega)\le M_i^{\max},\quad \forall i,\omega.\notag
\end{align}
Таким образом, оптимизация направлена на среднюю прогнозную \emph{задержку одного запроса} по заданной нагрузке, а не на несуществующее в исходном прототипе перекрытие микропакетов. Если ограничения памяти нарушаются, соответствующее разбиение недопустимо.

\subsection{Метод решения и проверка оптимальности в модели}

Поскольку порядок блоков фиксирован и стоимость~(\ref{eq:cost}) аддитивна по стадиям и границам, для модели предлагается динамическое программирование. Обозначим через $F(i,q)$ минимальную допустимую среднюю стоимость размещения первых $q$ блоков на первых $i$ исполнителях. Для каждой допустимой границы $p<q$:
\begin{equation}
\begin{split}
 F(i,q)=\min_{p}\bigl[&F(i-1,p)+\overline C_i(p,q)\\
 &+\overline Q_{i-1}(p)\bigr],
\end{split}
 \label{eq:dp}
\end{equation}
где $\overline C_i=\sum_\omega w_\omega C_i(\cdot;\omega)$ и аналогично $\overline Q_{i-1}$, а варианты с $M_i(p,q;\omega)>M_i^{\max}$ исключаются. Для первого исполнителя $Q_0=0$; $F(0,0)=0$, остальные недопустимые начальные состояния бесконечны.

\begin{proposition}
При фиксированном порядке узлов, аддитивной стоимости~(\ref{eq:cost}) и проверяемых для каждого сегмента ограничениях памяти рекуррентное соотношение~(\ref{eq:dp}) находит минимальную стоимость среди всех допустимых непрерывных разбиений модели.
\end{proposition}
\noindent\emph{Обоснование.} В любом оптимальном разбиении последняя стадия покрывает некоторый отрезок $[p,q)$. Если размещение предшествующих $p$ блоков не оптимально для первых $i-1$ исполнителей, его можно заменить более дешёвым допустимым размещением, уменьшив общую аддитивную стоимость. Поэтому оптимум обязан содержать оптимальную подзадачу и учитывается перебором всех $p$. По индукции получаем рекуррентное соотношение.

При заранее подготовленных стоимостях сегментов вычислительная сложность такого алгоритма составляет $O(KN^2)$, память для таблицы состояний~--- $O(KN)$. Для небольших $N$ корректность результата планируется проверять независимым полным перебором границ: например, при $N=6$, $K=3$ существует $\binom{5}{2}=10$ способов разделить блоки на три непустые последовательные группы.

Оптимальность полученного решения гарантируется только \emph{в рамках заданной модели затрат}. На MacBook, где три контейнера конкурируют за общие физические ядра и память, суммарное время может не быть аддитивным. Поэтому прогноз модели необходимо независимо сравнивать с фактическими замерами, а при сильном расхождении~--- пересматривать профилирование и ограничения. Для многозапросного перекрывающегося конвейера целевая функция будет другой; такая функциональность относится к дальнейшему развитию, а не к выполненному прототипу \cite{hwang2026}.

\subsection{Проект авторского набора данных и проверок}

Предусматривается собственный воспроизводимый набор из 144 нагрузочных сценариев: четыре категории текстов (продолжение фразы, вопросы, техническое описание, инструкции), три исходных текста в категории, четыре целевые длины входа (16, 64, 128, 256 токенов) и три длины генерации (8, 16, 32 токена). Число конфигураций равно $4\cdot3\cdot4\cdot3=144$. Тексты и правила их подготовки будут сохранены в файлах данных и генераторе; фактические длины определяются GPT-2-токенизатором, а не приписываются исходным текстам по целевому значению. Для каждой записи планируется фиксировать \code{case\_id}, \code{prompt\_text}, \code{input\_tokens}, \code{max\_new\_tokens}, версию модели и происхождение текста.

Набор не предназначен для обучения модели: это \emph{авторские входные нагрузки для измерения свойств системы}. Отдельно предполагается сохранять машиночитаемый журнал измерений по всем режимам, в том числе ошибочные запуски и несовпадения токенов. Исследовательский датасет и результаты измерений на момент ПЗ-3 ещё не созданы; здесь представлена проверяемая спецификация их получения.

\subsection{Экспериментальная методика в пределах доступных ресурсов}

Планируемая базовая аппаратная платформа~--- MacBook Pro с Apple M1 Pro, 16~ГБ объединённой памяти и приблизительно 146~ГБ свободного пространства. Docker Desktop на Mac использует Linux-виртуальную машину с отдельным лимитом оперативной памяти, по умолчанию 50\% памяти хоста \cite{dockermemory}. Следовательно, реальная доступность ресурсов для Ray определяется настройками Docker и нагрузкой macOS. В Linux-контейнерах нельзя автоматически рассчитывать на ускорение PyTorch через macOS MPS; основная серия задумывается как CPU-эксперимент.

Развитие модели выполняется ступенчато: сначала воспроизводятся базовые тесты на tiny-gpt2, далее проверяется \code{distilbert/distilgpt2} (приблизительно 82 млн параметров по карточке модели), затем при достаточном запасе памяти~--- \code{openai-community/gpt2} (124 млн параметров) \cite{distilgpt2card,gpt2card}. При переходе необходимо избежать пикового дублирования полных весов, которое может возникнуть, если каждый Ray Actor сначала загружает всю модель, а затем удаляет ненужные части.

Для каждого этапа применяется последовательность проверок: \emph{(1)} запуск и проверка памяти на небольшой пилотной выборке; \emph{(2)} точная проверка идентификаторов токенов и диапазонов слоёв; \emph{(3)} повторные измерения после прогрева; \emph{(4)} расширение нагрузки только при успешных проверках. Времена сравниваются на одной физической машине, с фиксированными версиями библиотек, числом потоков PyTorch и ограничениями Ray Actors. Журналы должны содержать фактические значения CPU, RAM, версию окружения и ошибки. Запуски с нехваткой памяти не исключаются незаметно из анализа.

\begin{table*}[t]
\centering\small
\caption{Этапность и критерии допустимости исследований}
\label{tab:plan}
\begin{tabularx}{\textwidth}{@{}p{3.2cm}X p{3.0cm}@{}}
\toprule
\textbf{Этап} & \textbf{Критерий перехода} & \textbf{Статус}\\
\midrule
Исходный tiny-gpt2 / Ray & Успешная генерация, размещение трёх частей, базовые тесты & Реализовано в НИР\\
Четыре варианта логитов & Согласованные token ID, учёт вычислений и обмена & Планируется\\
Авторские 144 сценария & Проверенные токены, происхождение, версии и журналы & Планируется\\
DistilGPT2, затем GPT-2 & Устойчивый CPU-запуск без переполнения памяти & Условно планируется\\
Оптимизация разбиения & Сверка DP с перебором и независимым замером & Планируется\\
KV-кэш, многозапросный pipeline & Отдельная архитектурная и экспериментальная работа & Перспектива\\
Несколько физических серверов & Наличие дополнительного оборудования и сетевых замеров & Вне текущего стенда\\
\bottomrule
\end{tabularx}
\end{table*}

На стадии сравнения необходимо использовать одинаковые сценарии для всех режимов, фиксировать прогрев, повторять запуски и оценивать медиану и разброс задержек. Измерения полезного объёма тензоров следует отделять от инструментального измерения сетевых байтов и сериализации. Для моделей с одинаковой математикой допустимо ожидать совпадения токенов в детерминированном режиме, но факт совпадения должен быть подтверждён проверкой, а не объявлен заранее.

\subsection{Воспроизводимость и границы применимости}

Результирующий экспериментальный проект должен включать исходники, фиксированные зависимости, конфигурации запуска, параметры аппаратной среды, генератор нагрузочного датасета, идентификаторы моделей, сохранённые метрики и скрипты построения графиков. Теоретические оценки, прогнозы модели распределения и измеренные показатели должны маркироваться раздельно.

Следует различать: (а) доказанное корректное выполнение учебного распределённого прототипа; (б) математическую возможность сократить передаваемое сообщение; (в) будущую проверку фактического ускорения; (г) дальнейшую задачу оптимального распределения слоёв. Уменьшение объёма данных не равно установленному ускорению, а оптимальность в аддитивной модели не равна оптимальности на реальном оборудовании. С учётом этих ограничений предлагаемый план остаётся реализуемым на доступной рабочей станции без обязательного использования GPU-кластера.

% Список литературы в формате оригинального шаблона MISIS.
% Библиография нумеруется автоматически в порядке первого упоминания.
\begin{thebibliography}{99}

\bibitem{petals2023}
Borzunov A., Baranchuk D., Dettmers T. et al. Petals: Collaborative Inference and Fine-tuning of Large Models // Proceedings of ACL: System Demonstrations. 2023. P. 558--568. DOI: \href{https://doi.org/10.18653/v1/2023.acl-demo.54}{10.18653/v1/2023.acl-demo.54}.

\bibitem{deepspeed2022}
Aminabadi R. Y., Rajbhandari S., Zhang M. et al. DeepSpeed Inference: Enabling Efficient Inference of Transformer Models at Unprecedented Scale. 2022. arXiv:2207.00032. \url{https://arxiv.org/abs/2207.00032}.

\bibitem{jekel2026}
Jekel C. F., Bhatele A. Understanding and Improving Communication Performance in Multi-node LLM Inference // Proceedings of the ACM Conference on AI and Agentic Systems. 2026. P. 1070--1083. DOI: \href{https://doi.org/10.1145/3786335.3813165}{10.1145/3786335.3813165}.

\bibitem{hwang2026}
Hwang S., Ahn J. Revisiting Pipeline Parallelism for LLM Serving // OSDI 2026. P. 1803--1819. \url{https://www.usenix.org/conference/osdi26/presentation/hwang}.

\bibitem{flexgen2023}
Sheng Y., Zheng L., Yuan B. et al. FlexGen: High-Throughput Generative Inference of Large Language Models with a Single GPU // Proceedings of ICML. 2023. Vol. 202. P. 31094--31116. \url{https://proceedings.mlr.press/v202/sheng23a.html}.

\bibitem{vllm2023}
Kwon W., Li Z., Zhuang S. et al. Efficient Memory Management for Large Language Model Serving with PagedAttention // Proceedings of SOSP. 2023. P. 611--626. DOI: \href{https://doi.org/10.1145/3600006.3613165}{10.1145/3600006.3613165}.

\bibitem{distserve2024}
Zhong Y., Liu S., Chen J. et al. DistServe: Disaggregating Prefill and Decoding for Goodput-optimized Large Language Model Serving // OSDI 2024. P. 193--210. \url{https://www.usenix.org/conference/osdi24/presentation/zhong-yinmin}.

\bibitem{ray2018}
Moritz P., Nishihara R., Wang S. et al. Ray: A Distributed Framework for Emerging AI Applications // OSDI 2018. P. 561--577. \url{https://www.usenix.org/conference/osdi18/presentation/moritz}.

\bibitem{rayresources}
Ray Team. Resources: Physical and Logical Resources // Ray Documentation. 2026. \url{https://docs.ray.io/en/latest/ray-core/scheduling/resources.html}.

\bibitem{dockermemory}
Docker Inc. Change your Docker Desktop settings: Resources // Docker Documentation. 2026. \url{https://docs.docker.com/desktop/settings-and-maintenance/settings/}.

\bibitem{distilgpt2card}
Hugging Face. DistilGPT2: Model Card. \url{https://huggingface.co/distilbert/distilgpt2}.

\bibitem{gpt2card}
OpenAI Community. GPT-2: Model Card. \url{https://huggingface.co/openai-community/gpt2}.

\end{thebibliography}

\end{document}
